\section{Playing by the Book}

\begin{quotex}
For my part I consider that it is better to be adventurous than cautious, because fortune is a woman, and if you wish to keep her under it is necessary to beat and ill-use her; and it is seen that she allows herself to be mastered by the adventurous rather than by those who go to work more coldly. She is, therefore, always, woman-like, a lover of young men, because they are less cautious, more violent, and with more audacity command her.
\flright{\textsc{Nicolo Machiavelli}, \emph{The Prince}}
\end{quotex}

Many years ago this writer used to play a lot of Backgammon. While it was a fad, local discos would host backgammon tournaments, with a small entrance fee. Since I was quite competent at the game, and the competition was generally weak, it was an easy way for me to pick up the \$50 or \$100 prize money without much effort at a time when that amount of money was meaningful to a young man.

At one particular tournament, with a larger prize, the competition was stiffer. I found myself paired against an older woman, typical of South Florida, who was clearly schooled in the game. She probably had taken lessons at her condo rec room. Finding myself in an untenable position, I deliberately took a calculated risk, a move that, if made by a patzer, would have been a sign of incompetence. I could sense the rush of excitement in my opponent as she anticipated moving on to the next table. Instead, my next throw was double sixes which knocked her out of the tournament.

In a voice thick with contempt, she told me, ``you know that move was incorrect'' and glared at me as though I had cheated. True enough, if you go by the backgammon manuals, but on the other hand it was the only move that gave me the opportunity to take the game. Had I ``played by the book'', I would have gone home early that night.


The world is full of those who would ``write the book'' to constrain Fortuna\footnote{\url{https://tinyurl.com/yc2wce35}}. The book will tell you how to behave, it will make false promises to those who follow the book, but the net result is to protect those in the positions described in the book from random or unexpected results. This is an attempt to control the future by limiting possibilities.

However, the metaphysical principle that ``the possible is the real'' will doom any such project to ultimate failure. Instead, the adept will understand the future by understanding the possibilities inherent in the moment. Evola writes in The Hermetic Tradition: ``\emph{prophetic knowledge, rather than being based on fate, is based on its mastery}'' and \textbf{Zosimos}: ``\emph{the race of Philosophers is beyond destiny.}'' \textbf{Plotinus} writes:

\begin{quotex}
The knowledge that the superior man has of the future that we attribute to him does not at all resemble the knowledge of fortune-tellers, but is like that of active participants who have a certainty of what must be; and such is the case with genuine rulers. For them nothing is undetermined, nothing is uncertain. Their decision persists as it was in the first moment. Their judgment of things to come is just as firm as their vision of present things is exact.
\end{quotex}


We \textbf{Sons of Hermes} understand possibilities, and master them by an act of will. Evola writes:

\begin{quotex}
Homer said that the gods often travel through the world in the guise of strangers and pilgrims and turn the cities of men around. This is not just mythology. There is reason to believe that no historical or social event of any importance, no phenomenon that has followed a determined course of terrestrial events comprising certain discoveries or the birth of new ideas, has not had a causal or spontaneous origin, but on the contrary has obeyed an \emph{intention}, if not an actual plan conceived behind the scenes and realized via paths we can scarcely imagine today, under the sign of the Light, as well as---according to circumstances---under the sign of Darkness.
\end{quotex}



\flrightit{Posted on 2011-06-01 by Cologero}

\begin{center}* * *\end{center}

\begin{footnotesize}
\begin{sffamily}
\texttt{Matt on 2017-06-02 at 11:54 said:}

This reminds me of a statement made by a character in the old X-Files series. To paraphrase: The best way to predict the future is to create it.

\end{sffamily}
\end{footnotesize}

